% context: markscheme (official solutions) for 2-3HW_Acceleration
% topic: simple calculations at constant acceleration with the four kinematic
%        equations; speeding up from rest, braking to a stop, a negative acceleration
% convention: answers in bold; "=" joins exact equalities, "\approx" introduces a
%        rounded value; each solution lists the knowns, then names the equation used
\documentclass[12pt, twoside]{article}
%\documentclass[12pt, twoside]{article}
\usepackage[letterpaper, margin=1in, headsep=0.2in]{geometry}
\setlength{\headheight}{0.6in}
%\usepackage[english]{babel}
\usepackage[utf8]{inputenc}
\usepackage{microtype}
\usepackage{amsmath}
\usepackage{amssymb}
%\usepackage{amsfonts}
\usepackage[nomessages]{fp} %\FPeval{\var-name}{2*sin(pi/6)}
\usepackage{siunitx} %units in math. eg 20\milli\meter
\usepackage{yhmath} % for arcs, overparenth command
\usepackage{tikz} %graphics
\usetikzlibrary{quotes, angles, arrows, arrows.meta}
\usepackage{graphicx} %consider setting \graphicspath{{images/}}
\usepackage{parskip} %no paragraph indent
\usepackage{enumitem}
\usepackage{multicol}
\usepackage{venndiagram}

\usepackage{fancyhdr}
\pagestyle{fancy}
\fancyhf{}
\renewcommand{\headrulewidth}{0pt} % disable the underline of the header
\raggedbottom
\hfuzz=2mm %suppresses overfull box warnings

\usepackage{hyperref}
\usepackage{float}
\setlength{\parskip}{0.2\baselineskip}

\title{Physics}
\author{Chris Huson}
\date{October 2026}

\fancyhead[LE]{\thepage}
\fancyhead[RO]{\thepage}
\fancyhead[LO]{La Scuola d'Italia / Huson / Physics \\* 23 October 2026}

\begin{document}

\subsubsection*{2.3 Homework: Constant Acceleration Solutions}

\begin{enumerate}[itemsep=0.5cm]

\item Vernier cart. Knowns: $v_0 = 0$, $a = 0.500$ m/s$^2$, $x_0 = 0$.
  \begin{enumerate}[itemsep=0.3cm]
    \item \renewcommand{\arraystretch}{1.4}
      \begin{tabular}[t]{| c | c | c | c | c | c |}
        \hline
        $t$ (s) & 0 & 1.00 & 2.00 & 3.00 & 4.00 \\ \hline
        $v$ (m/s) & \textbf{0} & \textbf{0.500} & \textbf{1.00} & \textbf{1.50}
          & \textbf{2.00} \\ \hline
      \end{tabular}

      \emph{Each second adds 0.500 m/s; this is $v = v_0 + at$ one row at a time.
      The 0.500 m/s at $t = 1.00$ s is the instantaneous velocity found in 2.2,
      question 4(b).}
    \item Want $x - x_0$ at $t = 2.00$ s. Use $x = x_0 + v_0 t + \frac{1}{2}at^2$:

      $x - x_0 = 0 + \tfrac{1}{2}(0.500~\text{m/s}^2)(2.00~\text{s})^2
      = \mathbf{1.00~\textbf{m}}$. \quad \textbf{Yes, it agrees} with the measured
      1.000 m.

      \emph{Also accept the average velocity route:
      $\bar{v} = \frac{1.00 + 0}{2} = 0.500$ m/s, and
      $0.500~\text{m/s} \times 2.00~\text{s} = 1.00$ m. The 2.2 data were made from
      $x = 0.25\,t^2$, which is exactly $\frac{1}{2}at^2$ with $a = 0.500$ m/s$^2$.}
  \end{enumerate}

\item Tesla Model S Plaid. Knowns: $v_0 = 0$, $v = 26.8$ m/s, $t = 1.99$ s.
  \begin{enumerate}[itemsep=0.3cm]
    \item Want $a$. Use $v = v_0 + at$, so $a = \dfrac{v - v_0}{t}$:

      $a = \dfrac{26.8~\text{m/s} - 0}{1.99~\text{s}} = 13.467\ldots
      \approx \mathbf{13.5~\textbf{m/s}^2}$

      \emph{The car gains 13.5 m/s of speed every second, more than a falling
      object's 9.80 m/s$^2$. Tesla's figure leaves out the first foot of travel, as
      drag-strip timing does, so the true average acceleration from a standstill is
      a little lower.}
    \item Want $x - x_0$. Use $\bar{v} = \dfrac{v + v_0}{2}$:
      $\bar{v} = \dfrac{26.8 + 0}{2} = 13.4$ m/s.

      $x - x_0 = \bar{v}\,t = 13.4~\text{m/s} \times 1.99~\text{s} = 26.666
      \approx \mathbf{26.7~\textbf{m}}$

      \emph{Also accept $x - x_0 = \frac{1}{2}at^2
      = \frac{1}{2}(13.467\ldots)(1.99)^2 \approx 26.7$ m. Using the rounded 13.5
      m/s$^2$ gives 26.73 m, which also rounds to 26.7 m.}
  \end{enumerate}

\newpage

\item Braking. Knowns: $v_0 = 13.4$ m/s, $v = 0$, $x - x_0 = 14$ m. Take the
  direction of travel as positive.
  \begin{enumerate}[itemsep=0.3cm]
    \item Want $a$; no time is given. Use $v^2 = v_0^2 + 2a(x - x_0)$:

      $0 = (13.4~\text{m/s})^2 + 2a(14~\text{m})$, \quad so

      $a = -\dfrac{(13.4~\text{m/s})^2}{2(14~\text{m})} = -\dfrac{179.56}{28}
      = -6.4128\ldots \approx \mathbf{-6.4~\textbf{m/s}^2}$

      \textbf{It is negative because the car is slowing down}: the acceleration points
      backward, opposite to the forward velocity.

      \emph{Two significant figures, because 14 m has two. Accept $-6.41$ m/s$^2$. A
      positive 6.4 m/s$^2$ described as ``deceleration'' earns the number but not the
      sign question.}
    \item Knowns: $v_0 = 26.8$ m/s, $v = 0$, $a = -6.4128\ldots$ m/s$^2$. Want
      $x - x_0$. Use $v^2 = v_0^2 + 2a(x - x_0)$ again:

      $x - x_0 = \dfrac{v^2 - v_0^2}{2a} = \dfrac{0 - (26.8~\text{m/s})^2}{2(-6.4128\ldots
      ~\text{m/s}^2)} = \dfrac{718.24}{12.8257\ldots} = \mathbf{56~\textbf{m}}$

      \textbf{Twice the speed gives four times the braking distance}: $4 \times 14$ m
      $= 56$ m.

      \emph{The distance grows with $v_0^2$, and $26.8 = 2 \times 13.4$, so
      $v_0^2$ is four times as large. The Highway Code prints 55 m; its table is
      rounded, and with two significant figures 55 m and 56 m agree. Accept 56.1 m
      from the rounded $-6.4$ m/s$^2$.}
  \end{enumerate}

\item Shanghai Tower elevator. Knowns: $v_0 = 0$, $v = 20.5$ m/s,
  $a = 1.0$ m/s$^2$.
  \begin{enumerate}[itemsep=0.3cm]
    \item Want $t$. Use $v = v_0 + at$, so $t = \dfrac{v - v_0}{a}$:

      $t = \dfrac{20.5~\text{m/s} - 0}{1.0~\text{m/s}^2} = 20.5
      \approx \mathbf{21~\textbf{s}}$

      \emph{Accept 20.5 s. The acceleration has two significant figures, so 21 s is
      the strict answer; either earns the mark.}
    \item Want $x - x_0$. Use $\bar{v} = \dfrac{v + v_0}{2} = \dfrac{20.5 + 0}{2}
      = 10.25$ m/s:

      $x - x_0 = \bar{v}\,t = 10.25~\text{m/s} \times 20.5~\text{s} = 210.125
      \approx \mathbf{210~\textbf{m}}$

      \emph{Also accept $v^2 = v_0^2 + 2a(x - x_0)$:
      $x - x_0 = \frac{(20.5)^2}{2(1.0)} = 210.125$ m, or
      $\frac{1}{2}at^2 = \frac{1}{2}(1.0)(20.5)^2$. Carry 20.5 s, not the rounded
      21 s, into this step.}
    \item Speeding up takes 210 m and slowing down takes another 210 m:
      $2 \times 210.125~\text{m} = 420.25 \approx \mathbf{420~\textbf{m}}$.

      \emph{On any ride shorter than about 420 m the elevator must start slowing down
      before it reaches 20.5 m/s. The tower is 632 m tall, so a long ride such as
      the express from the lobby to the observation deck gets there; short hops
      between floors never do.}
  \end{enumerate}

\end{enumerate}
\end{document}
